\section{The Frozen Bank}
\label{sec:results}

\subsection{Bank composition and quality}
Quality control reduces the $1{,}823$ input items to a frozen healthy
bank of \textbf{1{,}284} items (Table~\ref{tab:qc}). The four drop
reasons account for 539 removed items; zero-variance and non-positive
point-biserial dominate, with relatively few items failing the fit
bounds.

\begin{table}[t]
\centering
\small
\begin{tabular}{@{}lr@{}}
\toprule
stage / reason & items \\
\midrule
input items                 & $1{,}823$ \\
\quad dropped: zero-variance        & 308 \\
\quad dropped: point-biserial $\le 0$ & 155 \\
\quad dropped: infit $> 1.5$        & 14 \\
\quad dropped: outfit $> 2.0$       & 62 \\
\midrule
\textbf{frozen (healthy)}   & \textbf{1{,}284} \\
\bottomrule
\end{tabular}
\caption{Phase~2 quality-control funnel. Retention $\approx 70\%$.}
\label{tab:qc}
\end{table}

The healthy bank has median difficulty $b = +1.65$, median point-biserial
$+0.42$, median infit $0.87$, and median outfit $0.66$
(Table~\ref{tab:quality}). Difficulty spans roughly $\pm 4$ logits
(Figure~\ref{fig:bhist}); the item-fit cloud sits comfortably inside the
QC bounds (Figure~\ref{fig:itemfit}).

\begin{table}[t]
\centering
\small
\begin{tabular}{@{}lcccc@{}}
\toprule
 & $b$ & point-biserial & infit & outfit \\
\midrule
median (healthy bank) & $+1.65$ & $+0.42$ & $0.87$ & $0.66$ \\
\bottomrule
\end{tabular}
\caption{Healthy-bank summary statistics ($n = 1{,}284$).}
\label{tab:quality}
\end{table}

\paragraph{Difficulty with uncertainty.}
We report difficulty as $b \pm \mathrm{SE}$. The standard error shown
here is the analytic delta-method value for a single-item logit
difficulty,
$\mathrm{SE}(b) = 1/\sqrt{N\,P(1-P)}$ with $N = 45$ and
$P = \sigma(-b)$. This is an interim estimate; the posterior standard
deviation from a variational Rasch fit (\texttt{py-irt}) is the intended
final uncertainty and is wired into the pipeline
(\texttt{evaluator/pyirt\_fit.py}) but requires regenerating the raw
response matrix; we report the analytic value here and treat the
posterior comparison as future work (Section~\ref{sec:discussion}).

\begin{figure}[t]
\centering
\includegraphics[width=0.78\linewidth]{fig1_b_hist.pdf}
\caption{Difficulty distribution of the frozen bank ($n = 1{,}284$),
baseline-anchored to $\theta = 0$. Median $b = +1.65$.}
\label{fig:bhist}
\end{figure}

\begin{figure}[t]
\centering
\includegraphics[width=0.62\linewidth]{fig2_item_fit.pdf}
\caption{Item fit (infit vs.\ outfit) for the frozen bank. Dashed lines
are the QC upper bounds (infit $1.5$, outfit $2.0$); dotted lines mark
the ideal Rasch value of $1.0$.}
\label{fig:itemfit}
\end{figure}

\subsection{Real-vs-synthetic difficulty gap}
Items extracted from FDA warning letters are systematically harder than
synthetic items on the frozen bank: $n = 93$ real items with mean
$b = +1.98$ versus $n = 1{,}191$ synthetic with mean $b = +1.04$, a gap
of $+0.94$ logits. A tie-corrected Mann--Whitney $U$ test confirms the
shift is not attributable to chance ($U = 66{,}516$, $z = -3.24$,
two-sided $p = 0.0012$). The effect is real but modest in magnitude:
Cliff's $\delta = +0.20$, i.e.\ a randomly drawn real item is harder than
a randomly drawn synthetic one about $60\%$ of the time. The difference
is better described as a shift in the \emph{lower tail} than a wholesale
separation: only $5$ of $93$ real items ($5.4\%$) fall below $b = 0$,
against $372$ of $1{,}191$ synthetic items ($31.2\%$), and the easiest
real item ($b = -2.00$) is still far above the easiest synthetic one
($b = -4.25$). Enforcement, in other words, rarely produces easy
questions --- but it does not produce exclusively hard ones either.

On the full Phase~1 pre-QC set ($n = 105$ real, $n = 1{,}718$ synthetic)
the mean gap is $+0.49$ logits, suggesting QC preferentially removes
easy-or-inverted synthetic items and thus sharpens the real-vs-synthetic
separation. The test is reproducible from the committed frozen bank with
no API calls via \texttt{scripts/real\_vs\_synthetic\_test.py}.
